% Part: normal-modal-logic
% Chapter: axioms-systems
% Section: duals

\documentclass[../../../include/open-logic-section]{subfiles}

\begin{document}

\olfileid[fr]{nml}{prf}{dua}

\olsection{\usetoken{P}{formula} duales}

\begin{defn}\ollabel{def:duals}
  Chacune des !!{formula}s \Ax{T}, \Ax{B}, \Ax{4} et
  \Ax{5} possède une \emph{duale}, notée par un losange en indice,
  comme suit :
  \begin{align}
    \tag{\Ax{T_\Diamond}} p & \lif \Diamond p\\
    \tag{\Ax{B_\Diamond}} \Diamond\Box p & \lif p\\
    \tag{\Ax{4_\Diamond}} \Diamond\Diamond p & \lif \Diamond p\\
    \tag{\Ax{5_\Diamond}} \Diamond\Box p & \lif \Box p
    \end{align}
\end{defn}

Chacune des !!{formula}s duales ci-dessus s'obtient à partir de la
!!{formula} correspondante en substituant $\lnot p$ à $p$, en
prenant la contraposée, en remplaçant $\lnot\Box\lnot$ par $\Diamond$
et en remplaçant $\lnot\Diamond\lnot$ par~$\Box$. La formule \Ax{D},
c'est-à-dire $\Box!A \lif \Diamond!A$, est sa propre duale en ce sens.

\begin{prop}\ollabel{prop:dualsys}
  Pour chaque !!{formula}~$!A$ de la \olref[nml][prf][dua]{def:duals} :
  $\Log{K}!A = \Log{K}!A_{\Diamond}$.
\end{prop}
\begin{proof}
  Exercice.
\end{proof}
\begin{prob}
  Démontrer la \olref[nml][prf][dua]{prop:dualsys}.
\end{prob}

\end{document}
